% संपूर्ण मराठी ग्रंथातील प्रकरण 77: ग्रीक वर्णमाला
% हा संपादनयोग्य स्रोतखंड आहे. संकलनासाठी संपूर्ण स्रोतसंग्रहातील
% अक्षररूपे, पूर्वभाग, चिन्हव्याख्या आणि आकृती-स्रोत आवश्यक आहेत.
% मूळ: Open Logic Project; मजकूर आणि रूपांतर CC BY 4.0.
% यंत्रानुवाद, दुरुस्ती आणि तपासणी: OpenAI Codex —
% GPT-5.6 Sol आणि GPT-6 Sol, दोन्ही Ultra effort.
% दुरुस्ती-पुनरावलोकन: GPT-6.1 Sol, Ultra effort.
\chapter{ग्रीक वर्णमाला}\label{ref:grk:chap}








\begin{center}
    \begin{tabular}{lll@{\qquad}lll}
    अल्फा &	$\alpha$ & $A$ &
    न्यू & $\nu$ & $N$ \\
    बीटा & $\beta$ & $B$ &
    झाय & $\xi$ & $\Xi$ \\
    गॅमा &	$\gamma$ & $\Gamma$ &
    ओमिक्रॉन & $o$ & $O$ \\
    डेल्टा & $\delta$ & $\Delta$ &
    पाय & $\pi$ & $\Pi$ \\
    एप्सिलॉन & $\epsilon$ & $E$ &
    रो & $\rho$ & $P$ \\
    झीटा & $\zeta$ & $Z$ &
    सिग्मा &	$\sigma$ & $\Sigma$ \\
    एटा & $\eta$ & $H$ &
    टाऊ & $\tau$ & $T$ \\
    थीटा & $\theta$ & $\Theta$ &
    अप्सिलॉन & $\upsilon$ & $\Upsilon$ \\
    आयोटा & $\iota$ & $I$ &
    फाय & $\varphi$ & $\Phi$ \\
    कॅप्पा & $\kappa$ & $K$ &
    खाय & $\chi$ & $X$ \\
    लॅम्ब्डा & $\lambda$ & $\Lambda$ &
    साय & $\psi$ & $\Psi$ \\
    म्यू & $\mu$ & $M$ &
    ओमेगा & $\omega$ & $\Omega$
    \end{tabular}
\end{center}



